% ==========================================================================
%  Prova 1 — Gabarito Comentado
%  Macroeconomia I (PPGEco-UFES, 2025)
%  Prof. Dr. Ricardo Ramalhete Moreira
%  Cobre: RCK (cap.2), RBC (cap.5), Rigidez Nominal (cap.6),
%         Lucas / Fischer / Taylor / Calvo (cap.7)
% ==========================================================================
\documentclass[12pt, a4paper]{article}

\usepackage[utf8]{inputenc}
\usepackage[T1]{fontenc}
\usepackage{lmodern}
\usepackage[brazil]{babel}
\usepackage{microtype}

\usepackage[top=2.5cm, bottom=2.5cm, left=3cm, right=3cm, headheight=15pt]{geometry}

\usepackage{amsmath, amssymb, amsthm}
\usepackage{mathtools}
\usepackage{bm}

\usepackage{booktabs}
\usepackage{array}
\usepackage{multicol}
\usepackage{xcolor}
\usepackage{hyperref}
\usepackage{enumitem}
\usepackage{tcolorbox}
\tcbuselibrary{breakable}
\usepackage{fancyhdr}
\usepackage{titlesec}

% --------------------------------------------------------------------------
% Cores
% --------------------------------------------------------------------------
\definecolor{cyanrbc}{HTML}{3b9dbd}
\definecolor{orangerbc}{HTML}{d97b39}
\definecolor{grayrbc}{HTML}{6b7280}
\definecolor{greenbox}{HTML}{166534}
\definecolor{greenbg}{HTML}{dcfce7}
\definecolor{boxbg}{HTML}{f5f0e8}
\definecolor{boxborder}{HTML}{3b9dbd}
\definecolor{provabg}{HTML}{fff7ed}
\definecolor{provaborder}{HTML}{d97b39}
\definecolor{correctbg}{HTML}{dcfce7}
\definecolor{correctborder}{HTML}{166534}

% --------------------------------------------------------------------------
% Caixas
% --------------------------------------------------------------------------
\newtcolorbox{resultbox}[1][]{standard, colback=boxbg, colframe=boxborder,
  boxrule=1pt, arc=3pt, left=6pt, right=6pt, top=4pt, bottom=4pt,
  fonttitle=\bfseries\small, #1}
\newtcolorbox{provabox}[1][]{standard, colback=provabg, colframe=provaborder,
  boxrule=1.2pt, arc=3pt, left=6pt, right=6pt, top=4pt, bottom=4pt,
  fonttitle=\bfseries\small, #1}
\newtcolorbox{correctbox}[1][]{standard, colback=correctbg, colframe=correctborder,
  boxrule=1pt, arc=3pt, left=6pt, right=6pt, top=4pt, bottom=4pt,
  fonttitle=\bfseries\small\color{correctborder}, #1}
\newtcolorbox{formulabox}[1][]{standard, colback=greenbg!30, colframe=greenbox!60,
  boxrule=0.8pt, arc=3pt, left=6pt, right=6pt, top=4pt, bottom=4pt,
  fonttitle=\bfseries\small, #1}

% --------------------------------------------------------------------------
% Comandos matemáticos
% --------------------------------------------------------------------------
\newcommand{\E}{\mathbb{E}}
\newcommand{\dif}{\mathrm{d}}
\newcommand{\pibar}{\bar{\pi}}

\setlength{\parskip}{0.5\baselineskip}
\setlength{\parindent}{0pt}

\titleformat{\section}{\large\bfseries\color{cyanrbc}}{Questão \thesection.}{0.5em}{}[\vspace{-0.3em}\rule{\linewidth}{0.4pt}]
\titleformat{\subsection}{\normalsize\bfseries\color{grayrbc}}{}{0em}{}

\pagestyle{fancy}
\fancyhf{}
\lhead{\footnotesize\textcolor{grayrbc}{P1 2025 --- Gabarito Comentado --- Macroeconomia I (PPGEco-UFES)}}
\rhead{\footnotesize\textcolor{grayrbc}{RCK $\cdot$ RBC $\cdot$ Rigidez Nominal $\cdot$ NK}}
\cfoot{\small\thepage}
\renewcommand{\headrulewidth}{0.4pt}

% ==========================================================================
\begin{document}
% ==========================================================================

% --------------------------------------------------------------------------
% Capa
% --------------------------------------------------------------------------
\begin{titlepage}
\begin{center}
\vspace*{1.5cm}
{\Large\bfseries Universidade Federal do Espírito Santo}\\[0.5em]
{\large Programa de Pós-Graduação em Economia (PPGEco-UFES)}\\[2em]

\rule{\linewidth}{1.2pt}\\[0.8em]
{\LARGE\bfseries\color{cyanrbc} Prova 1 --- Gabarito Comentado\\[0.4em]
\large Macroeconomia I --- 2025}\\[0.4em]
{\normalsize\color{grayrbc} Romer caps.\ 2, 5, 6 e 7}\\[0.8em]
\rule{\linewidth}{1.2pt}

\vspace{1em}
{\large Prof.\ Dr.\ Ricardo Ramalhete Moreira}

\vspace{1.5em}
\begin{tcolorbox}[standard, colback=provabg, colframe=provaborder, width=0.88\textwidth,
  arc=4pt, boxrule=1.2pt, left=8pt, right=8pt, top=6pt, bottom=6pt]
\centering
\textbf{\textcolor{provaborder}{Conteúdo:}} 14 questões de múltipla escolha (escolha única)\\[0.4em]
\textbf{RCK} (Q1--Q4) $\cdot$ \textbf{RBC} (Q5--Q7) $\cdot$ \textbf{Rigidez Nominal} (Q8--Q9)\\[0.2em]
\textbf{Oferta de Trabalho} (Q10) $\cdot$
\textbf{Lucas / NK / Fischer / Taylor / Calvo} (Q11--Q14)
\end{tcolorbox}

\vspace{1.5em}
\begin{tcolorbox}[standard, colback=correctbg, colframe=correctborder,
  width=0.88\textwidth, arc=4pt, boxrule=1pt, left=8pt, right=8pt,
  top=6pt, bottom=6pt]
\centering
\textbf{\textcolor{correctborder}{Gabarito Rápido}}\\[0.6em]
\begin{tabular}{ccccccc}
\toprule
Q1 & Q2 & Q3 & Q4 & Q5 & Q6 & Q7 \\
\textbf{c} & \textbf{a} & \textbf{b} & \textbf{c} & \textbf{c} & \textbf{b} & \textbf{b} \\
\midrule
Q8 & Q9 & Q10 & Q11 & Q12 & Q13 & Q14 \\
\textbf{c} & \textbf{b} & \textbf{b} & \textbf{b} & \textbf{b} & \textbf{c} & \textbf{b} \\
\bottomrule
\end{tabular}
\end{tcolorbox}

\vfill
{\small Gabarito comentado elaborado com base na P1 de Macroeconomia I --- PPGEco-UFES 2025.}
\end{center}
\end{titlepage}

% ==========================================================================
% BLOCO RCK
% ==========================================================================
\newpage
\begin{center}
{\large\bfseries\color{cyanrbc} Bloco I --- Modelo Ramsey-Cass-Koopmans (Cap.\ 2)}\\[0.2em]
\rule{\linewidth}{0.6pt}
\end{center}

\setcounter{section}{0}

% ------------------------------------------------------------------
\section{Função Utilidade CRRA no Modelo RCK}
% ------------------------------------------------------------------

\textbf{Enunciado.} Seja a clássica função utilidade no modelo Ramsey-Cass-Koopmans (RCK):
\[
u(C) = \frac{C^{1-\theta}}{1-\theta},
\]
em que $\theta$ é o grau de aversão ao risco relativo. Identifique a afirmativa correta:

\begin{enumerate}[label=\alph*)]
\item O grau de aversão ao risco pode ser interpretado como a própria elasticidade de substituição intertemporal do consumo.
\item À medida que o consumo corrente aumenta, a utilidade marginal do consumo futuro vai se tornando cada vez menor.
\item A partir de tal função, deduz-se que o aumento do consumo corrente \emph{e} do grau de aversão ao risco reduzem a utilidade marginal.
\item A elevação do consumo tem efeito sobre a utilidade marginal $u'$, expresso pela forma $u' = -\theta\, C^{-\theta-1}$.
\end{enumerate}

\begin{correctbox}[title={Resposta correta: c)}]
A utilidade marginal é $u'(C) = C^{-\theta}$.

\textbf{Efeito de $C\uparrow$:}
\[
\frac{\partial u'}{\partial C} = -\theta\, C^{-\theta-1} < 0 \quad \forall\, C > 0.
\]
A MU é estritamente decrescente no consumo (utilidade marginal decrescente). \checkmark

\textbf{Efeito de $\theta\uparrow$:}
\[
\frac{\partial u'}{\partial\theta} = -\ln(C)\,C^{-\theta}.
\]
Para $C > 1$ (caso habitual no estado estacionário): $\ln C > 0$, portanto $\partial u'/\partial\theta < 0$ ---
maior aversão ao risco também reduz a MU. \checkmark

\textbf{Por que as demais são falsas:}
\begin{itemize}[leftmargin=*]
\item[a)] $\theta$ é o coeficiente de aversão relativa ao risco (CRRA); a elasticidade de
substituição intertemporal é $\mathrm{EIS} = 1/\theta$. São recíprocos, não iguais.
\item[b)] Com utilidade separável, $u'(C_{t+1})$ não depende diretamente de $C_t$; a conexão
só emerge via equação de Euler, não pela definição de $u'$.
\item[d)] O que é descrito é a \emph{segunda} derivada $u''(C) = -\theta C^{-\theta-1}$, não a primeira
derivada $u'(C) = C^{-\theta}$.
\end{itemize}
\end{correctbox}

% ------------------------------------------------------------------
\section{Equação de Euler para o Consumo no Modelo RCK}
% ------------------------------------------------------------------

\textbf{Enunciado.} Com base na equação de Euler para o consumo per capita no modelo RCK, identifique a correta:

\begin{enumerate}[label=\alph*)]
\item Em momentos em que o consumidor reage mais sensivelmente a mudanças na taxa real de juros, ceteris paribus, pode-se dizer que houve uma redução do grau de aversão ao risco.
\item A taxa de desconto afeta positivamente a taxa de crescimento do consumo.
\item Esta equação é deduzida a partir de uma função utilidade sujeita à restrição
$\int_0^\infty e^{-\rho t}C(t)L(t)\,\dif t \leq W^*$, onde $W^*$ é a riqueza da família.
\item A equação de Euler estabelece a necessidade de que famílias otimizadoras \emph{diminuam} o crescimento do consumo frente à elevação de juros reais.
\end{enumerate}

\begin{correctbox}[title={Resposta correta: a)}]
A Regra de Keynes-Ramsey:
\[
\frac{\dot{c}}{c} = \frac{f'(k) - \delta - \rho - \theta g}{\theta}
  = \frac{1}{\theta}\bigl(r - r^*\bigr),
  \qquad r^* = \rho + \theta g.
\]
A sensibilidade do crescimento do consumo à taxa de juros é $1/\theta$ (a EIS).
\textbf{Maior sensibilidade} $\Leftrightarrow$ maior $1/\theta$ $\Leftrightarrow$ \textbf{menor $\theta$}:
o consumidor reage mais a variações em $r$ quanto menor for o grau de aversão ao risco. \checkmark

\textbf{Por que as demais são falsas:}
\begin{itemize}[leftmargin=*]
\item[b)] $\rho\uparrow \Rightarrow r^*\uparrow \Rightarrow \dot{c}/c\downarrow$: a taxa de desconto afeta
\emph{negativamente} o crescimento do consumo.
\item[c)] A restrição correta é na forma fluxo: $\dot{k} = f(k) - c - (n+g+\delta)k$.
\item[d)] Juros reais $\uparrow \Rightarrow \dot{c}/c\uparrow$ (poupar é mais rentável $\Rightarrow$ consumo
cresce mais depressa): efeito \emph{positivo}, não negativo.
\end{itemize}
\end{correctbox}

% ------------------------------------------------------------------
\section{Equação de Movimento do Capital no Modelo RCK}
% ------------------------------------------------------------------

\textbf{Enunciado.} A dinâmica do capital por trabalho efetivo no modelo RCK depende basicamente da diferença entre investimento efetivo e desejado (\textit{breakeven investment}). Com base na equação de movimento do capital ($\dot{k}$), identifique a afirmativa correta:

\begin{enumerate}[label=\alph*)]
\item Uma elevação do crescimento do conhecimento ($g$), ceteris paribus, aumenta o estoque de capital na economia.
\item Quanto maior o produto marginal do capital, tudo o mais constante, menor precisa ser o consumo por trabalho efetivo a fim de manter estável o atual nível de $k$.
\item O crescimento populacional está positivamente relacionado com a variação do capital.
\item O crescimento equilibrado no modelo RCK ocorre na situação chamada de Regra de Ouro, quando o consumo por trabalho efetivo é maximizado.
\end{enumerate}

\begin{correctbox}[title={Resposta correta: b)}]
No locus $\dot{k}=0$: $c = f(k) - (n+g+\delta)k$.

Com Cobb-Douglas $f(k)=k^\alpha$, o consumo sustentável no locus é:
\[
c^*(k) = k^\alpha - (n+g+\delta)k.
\]
Derivando: $\dfrac{dc^*}{dk} = f'(k) - (n+g+\delta)$.

Abaixo da Regra de Ouro (caso normal): $f'(k) > n+g+\delta$, logo $dc^*/dk > 0$, ou seja,
$k$ menor implica $c^*$ menor. Como $f'(k) = \alpha k^{\alpha-1}$ é decrescente,
\textbf{maior $f'(k)$} corresponde a \textbf{menor $k$}, que por sua vez implica
\textbf{menor $c^*$} para manter $\dot{k}=0$. \checkmark

\textbf{Por que as demais são falsas:}
\begin{itemize}[leftmargin=*]
\item[a)] $g\uparrow$ eleva o \textit{breakeven} $(n+g+\delta)k$ $\Rightarrow$ $\dot{k}\downarrow$ $\Rightarrow$
$k^*\downarrow$. O estoque de capital \emph{diminui} no longo prazo.
\item[c)] $n\uparrow \Rightarrow (n+g+\delta)k\uparrow \Rightarrow \dot{k}\downarrow$: crescimento
populacional reduz (não aumenta) a variação do capital.
\item[d)] No RCK o longo prazo segue a \textbf{Regra de Ouro Modificada}:
$f'(k^*)=\rho+\theta g+\delta$, que em geral não maximiza $c^*$. A Regra de Ouro
($f'=n+g+\delta$, maximizadora de $c^*$) pertence ao modelo de Solow.
\end{itemize}
\end{correctbox}

% ------------------------------------------------------------------
\section{Diagrama de Fases do Modelo RCK}
% ------------------------------------------------------------------

\textbf{Enunciado.} Com base no diagrama de fases do modelo RCK, considere as seguintes situações e escolha apenas a verdadeira:

\begin{enumerate}[label=\alph*)]
\item Em uma situação à direita do locus $\dot{c}=0$ e acima do locus $\dot{k}=0$,
a economia aumenta o estoque de capital e reduz o consumo.
\item Em uma situação à esquerda do locus $\dot{c}=0$ e abaixo do locus $\dot{k}=0$,
o estoque de capital aumenta e há uma redução do consumo.
\item Em uma situação à direita do locus $\dot{c}=0$ e abaixo do locus $\dot{k}=0$,
existe uma elevação do capital e redução do consumo.
\item Em uma situação à esquerda do locus $\dot{c}=0$ e acima do locus $\dot{k}=0$,
a economia eleva o consumo e também o capital.
\end{enumerate}

\begin{correctbox}[title={Resposta correta: c)}]
Lógica do diagrama de fases (ponto de sela):

\begin{center}
\begin{tabular}{lll}
\toprule
Posição & Sinal & Efeito \\
\midrule
Direita do locus $\dot{c}=0$ ($k>k^*$) & $\dot{c}<0$ & $c$ cai \\
Esquerda do locus $\dot{c}=0$ ($k<k^*$) & $\dot{c}>0$ & $c$ sobe \\
Acima do locus $\dot{k}=0$ & $\dot{k}<0$ & $k$ cai \\
Abaixo do locus $\dot{k}=0$ & $\dot{k}>0$ & $k$ sobe \\
\bottomrule
\end{tabular}
\end{center}

Região c): \textbf{direita} ($k>k^*$) $\Rightarrow$ $\dot{c}<0$ ($c$ cai);
\textbf{abaixo} do $\dot{k}=0$ $\Rightarrow$ $\dot{k}>0$ ($k$ sobe).
Portanto: $k\uparrow$ e $c\downarrow$. \checkmark

\textbf{Análise das demais:}
\begin{itemize}[leftmargin=*]
\item[a)] Direita $+$ acima: $\dot{c}<0$ e $\dot{k}<0$ $\Rightarrow$ \emph{ambos} caem.
\item[b)] Esquerda $+$ abaixo: $\dot{c}>0$ e $\dot{k}>0$ $\Rightarrow$ ambos sobem; consumo não cai.
\item[d)] Esquerda $+$ acima: $\dot{c}>0$ e $\dot{k}<0$ $\Rightarrow$ $c$ sobe mas $k$ \emph{cai}.
\end{itemize}
\end{correctbox}

% ==========================================================================
% BLOCO RBC
% ==========================================================================
\newpage
\begin{center}
{\large\bfseries\color{cyanrbc} Bloco II --- Modelo de Ciclos Reais de Negócios (Cap.\ 5)}\\[0.2em]
\rule{\linewidth}{0.6pt}
\end{center}

% ------------------------------------------------------------------
\section{Substituição Intertemporal da Oferta de Trabalho}
% ------------------------------------------------------------------

\textbf{Enunciado.} No modelo Real Business Cycle (RBC) são considerados choques reais que estimulam variações no nível de atividade de curto prazo. A substituição intertemporal da oferta de trabalho é um fator importante no modelo. Escolha a alternativa correta:

\begin{enumerate}[label=\alph*)]
\item Um aumento de taxa real de juros estimula elevação do lazer no presente, ou seja, aumenta a oferta de trabalho futura.
\item A razão de substituição intertemporal pode ser expressa a partir de um problema de maximização em que a função utilidade do consumo é $u(C) = \dfrac{C^{1-\theta}-1}{1-\theta}$.
\item A reação de substituição intertemporal do trabalho é, em si, um fator explicativo do ciclo a partir de choques tecnológicos, já que estes últimos implicam em mudanças na taxa de juros e, consequentemente, da oferta de trabalho entre períodos.
\item A inclusão da incerteza no modelo inviabiliza mudanças na oferta de trabalho.
\end{enumerate}

\begin{correctbox}[title={Resposta correta: c)}]
No RBC, o mecanismo de propagação opera em três etapas:
\[
\text{Choque de PTF} \;\Rightarrow\; r\uparrow
  \;\Rightarrow\; \text{Substituição intertemporal de trabalho}
  \;\Rightarrow\; L_t\uparrow,\; L_{t+1}\downarrow.
\]
O choque tecnológico positivo eleva o produto marginal do capital (MPK), o que aumenta $r$.
A taxa de juros mais alta torna o trabalho \emph{hoje} mais lucrativo:
famílias ofertam mais trabalho no período corrente
(\textit{working while the sun shines}) e menos no futuro.
Esse mecanismo amplifica a resposta cíclica ao choque. \checkmark

\textbf{Por que as demais são falsas:}
\begin{itemize}[leftmargin=*]
\item[a)] $r\uparrow$ torna o lazer \emph{hoje} mais caro $\Rightarrow$ oferta de trabalho \emph{presente}
sobe e lazer migra para o futuro (não o contrário).
\item[b)] A afirmativa é verdadeira como descrição de uma forma funcional, mas não captura
o \emph{mecanismo} dinâmico que torna o RBC relevante para explicar o ciclo.
\item[d)] Incerteza gera comportamento de precaução e pode alterar poupança e trabalho;
não os inviabiliza.
\end{itemize}
\end{correctbox}

% ------------------------------------------------------------------
\section{RBC com Utilidade Log e Lazer: Condições de Primeira Ordem}
% ------------------------------------------------------------------

\textbf{Enunciado.} Seja o modelo RBC com a seguinte função utilidade:
\[
u = \ln c_1 + b\ln(1-l_1) + \beta\!\left[\ln c_2 + b\ln(1-l_2)\right], \quad b > 0,
\]
onde $c_i$ são os consumos e $l_i$ as ofertas de trabalho nos períodos $i=1,2$; $\beta$ é a taxa de desconto. A restrição orçamentária é:
\[
c_1 + \frac{c_2}{1+r} = w_1 l_1 + \frac{w_2 l_2}{1+r},
\]
sendo $w_i$ os salários reais e $r$ a taxa de juros. Identifique a questão correta:

\begin{enumerate}[label=\alph*)]
\item A expressão de Lagrange para a maximização é:
\[
\mathcal{L} = \ln c_1 + b\ln(1-l_1) + e^{-\rho}\!\left[\ln c_2 + b\ln(1-l_2)\right]
  + \lambda\!\left[w_1 l_1 + \frac{(1-l_2)}{1+r} - c_1 - \frac{c_2}{1+r}\right].
\]
\item Tomando as condições de primeira ordem para $l_1$ e $l_2$, chega-se à expressão
$\dfrac{1-l_2}{1-l_1} = \beta(1+r)\dfrac{w_1}{w_2}$, que descreve a razão de lazer intertemporal.
\item Pode-se afirmar que a utilidade marginal do trabalho é positiva.
\item Um processo de calibração do modelo seria, por exemplo, calcular todas as possíveis condições de primeira ordem para a função utilidade.
\end{enumerate}

\begin{correctbox}[title={Resposta correta: b)}]
\textbf{Lagrangiano correto:}
\[
\mathcal{L} = \ln c_1 + b\ln(1-l_1) + \beta[\ln c_2 + b\ln(1-l_2)]
  + \lambda\!\left[w_1 l_1 + \frac{w_2 l_2}{1+r} - c_1 - \frac{c_2}{1+r}\right].
\]

\textbf{CPOs:}
\begin{alignat}{2}
&\frac{\partial\mathcal{L}}{\partial c_1}=0: &\quad& \lambda = \frac{1}{c_1}. \\[4pt]
&\frac{\partial\mathcal{L}}{\partial l_1}=0: &\quad& \frac{b}{1-l_1} = \lambda w_1
  \;\Longrightarrow\; 1-l_1 = \frac{bc_1}{w_1}. \\[4pt]
&\frac{\partial\mathcal{L}}{\partial l_2}=0: &\quad& \frac{\beta b}{1-l_2} = \frac{\lambda w_2}{1+r}
  \;\Longrightarrow\; 1-l_2 = \frac{\beta b(1+r)c_1}{w_2}.
\end{alignat}

\textbf{Razão de lazer:}
\[
\frac{1-l_2}{1-l_1}
  = \frac{\beta b(1+r)c_1/w_2}{bc_1/w_1}
  = \beta(1+r)\frac{w_1}{w_2}. \quad\checkmark
\]
Interpretação: se $r\uparrow$ ou $w_1/w_2\uparrow$, o lazer migra para o período~2
(trabalha-se mais hoje).

\textbf{Por que as demais são falsas:}
\begin{itemize}[leftmargin=*]
\item[a)] O termo da restrição em a) usa $(1-l_2)/(1+r)$ em vez de $w_2 l_2/(1+r)$ ---
o Lagrangiano está incorreto.
\item[c)] O trabalho gera \emph{desutilidade}: $\partial u/\partial l_i = -b/(1-l_i) < 0$.
\item[d)] Calibração consiste em escolher parâmetros para replicar momentos dos dados,
não apenas derivar CPOs.
\end{itemize}
\end{correctbox}

% ------------------------------------------------------------------
\section{Efeitos de Choque Tecnológico Positivo no RBC}
% ------------------------------------------------------------------

\textbf{Enunciado.} Um choque tecnológico positivo no modelo RBC pode apresentar os seguintes efeitos. Escolha a alternativa correta:

\begin{enumerate}[label=\alph*)]
\item Uma vez que o choque vem acompanhado de elevação do consumo, há um estímulo à redução da poupança.
\item Há uma elevação do produto marginal do capital, o que impulsiona maior investimento das firmas, porém também o aumento de salários reais.
\item Embora haja uma elevação de salários reais, o mesmo não ocorre com a taxa de juros.
\item Quanto maior a persistência do choque tecnológico, a oferta de trabalho corrente tende a ficar abaixo do nível inicial.
\end{enumerate}

\begin{correctbox}[title={Resposta correta: b)}]
Um choque de PTF positivo ($A\uparrow$) com Cobb-Douglas $F = AK^\alpha L^{1-\alpha}$ eleva simultaneamente o MPK e o MPL:
\[
MPK = \alpha A k^{\alpha-1} \uparrow,
\qquad
MPL = (1-\alpha)A k^{\alpha} \uparrow.
\]
\begin{itemize}[leftmargin=*]
\item $MPK\uparrow$ $\Rightarrow$ rentabilidade do investimento sobe
  $\Rightarrow$ \textbf{investimento sobe}. \checkmark
\item $MPL\uparrow$ $\Rightarrow$ salário real de equilíbrio $w = MPL$
  \textbf{sobe}. \checkmark
\end{itemize}

\textbf{Por que as demais são falsas:}
\begin{itemize}[leftmargin=*]
\item[a)] Consumo \emph{e} poupança/investimento sobem após o choque;
a razão poupança/renda pode inclusive \emph{aumentar}.
\item[c)] $MPK\uparrow \Rightarrow r\uparrow$ também (mercado de capital em equilíbrio).
\item[d)] Maior persistência $\Rightarrow$ substituição intertemporal de trabalho ampliada
$\Rightarrow$ oferta de trabalho corrente \emph{sobe} (acima do nível inicial).
\end{itemize}
\end{correctbox}

% ==========================================================================
% BLOCO RIGIDEZ NOMINAL
% ==========================================================================
\newpage
\begin{center}
{\large\bfseries\color{cyanrbc}
  Bloco III --- Rigidez Nominal, Mark-up e Curva de Phillips (Cap.\ 6)}\\[0.2em]
\rule{\linewidth}{0.6pt}
\end{center}

% ------------------------------------------------------------------
\section{Rigidez Nominal, Mark-up e Ciclicidade dos Salários Reais}
% ------------------------------------------------------------------

\textbf{Enunciado.} Em modelos de rigidez nominal e competição imperfeita, a resposta do mercado de trabalho a choques de demanda depende basicamente do comportamento do \textit{mark-up} empresarial (assumindo-se rendimentos decrescentes do trabalho). Escolha a sentença verdadeira:

\begin{enumerate}[label=\alph*)]
\item Se os salários nominais e o mark-up são constantes, os salários reais tendem a aumentar diante de uma elevação da produção.
\item Se os salários nominais e o mark-up são constantes, um choque negativo de demanda vem acompanhado de queda dos salários reais.
\item Se o salário nominal é dado e o mark-up é contra-cíclico (cai mais do que o produto marginal do trabalho quando $Y$ sobe), os salários reais tornam-se pró-cíclicos.
\item O modelo keynesiano simples prevê que os salários reais são pró-cíclicos.
\end{enumerate}

\begin{correctbox}[title={Resposta correta: c)}]
O preço de equilíbrio com mark-up $\mu$ e desutilidade do trabalho é:
\[
P = \mu\,\frac{W}{\mathrm{MPL}} \;\Longrightarrow\; w \equiv \frac{W}{P} = \frac{\mathrm{MPL}}{\mu}.
\]
Com rendimentos decrescentes: $Y\uparrow \Rightarrow L\uparrow \Rightarrow \mathrm{MPL}\downarrow$,
o que sozinho \emph{reduziria} $w$.

Mas se $\mu$ é \textbf{contra-cíclico} e cai proporcionalmente mais que MPL durante uma expansão:
\[
\frac{\Delta w}{w} = \underbrace{\frac{\Delta\,\mathrm{MPL}}{\mathrm{MPL}}}_{<\,0}
  - \underbrace{\frac{\Delta\mu}{\mu}}_{\ll\,0} > 0
  \;\Longleftrightarrow\;
  \left|\frac{\Delta\mu}{\mu}\right| > \left|\frac{\Delta\,\mathrm{MPL}}{\mathrm{MPL}}\right|.
\]
Neste caso $w\uparrow$ quando $Y\uparrow$: \textbf{salários reais pró-cíclicos}. \checkmark

\textbf{Por que as demais são falsas:}
\begin{itemize}[leftmargin=*]
\item[a)] $\mu$ e $W$ constantes: $Y\uparrow \Rightarrow \mathrm{MPL}\downarrow \Rightarrow P\uparrow
  \Rightarrow w = W/P\downarrow$. Salários reais \emph{caem}.
\item[b)] Choque negativo: $Y\downarrow \Rightarrow \mathrm{MPL}\uparrow \Rightarrow P\downarrow
  \Rightarrow w\uparrow$. Salários reais \emph{sobem}.
\item[d)] No modelo keynesiano simples (preços rígidos, demanda determina produto),
$P$ não sobe com a demanda $\Rightarrow$ $w$ é \emph{acíclico} ou \emph{contra}cíclico,
não pró-cíclico.
\end{itemize}
\end{correctbox}

% ------------------------------------------------------------------
\section{Curva de Phillips Híbrida}
% ------------------------------------------------------------------

\textbf{Enunciado.} Seja uma Curva de Phillips dada por:
\[
\pi_t = \varphi\,\bar{\pi} + (1-\varphi)\,\pi_{t-1}
  + \lambda(\ln Y_t - \ln\bar{Y}) + \varepsilon_t.
\]
Indique a afirmativa correta:

\begin{enumerate}[label=\alph*)]
\item Uma situação próxima ao pleno emprego é mais propícia a maiores níveis de $\lambda$.
\item Sob desinflação da economia, a política monetária é beneficiada com graus menores para o componente inercial $(1-\varphi)$, ou seja, com $\varphi$ maior.
\item Esta equação pode ser deduzida diretamente da contribuição de Calvo (1983).
\item O componente $\varepsilon_t$ é relacionado à gestão monetária pelo Banco Central.
\end{enumerate}

\begin{correctbox}[title={Resposta correta: b)}]
Na equação híbrida:
\begin{itemize}[leftmargin=*]
\item $\varphi\bar{\pi}$: componente de ancoragem à meta (forward-looking / credibilidade do BC).
\item $(1-\varphi)\pi_{t-1}$: inércia inflacionária (backward-looking).
\end{itemize}

\textbf{Processo de desinflação.} O desvio da inflação em relação à meta satisfaz:
\[
\pi_t - \bar{\pi} = (1-\varphi)(\pi_{t-1} - \bar{\pi}) + \lambda x_t + \varepsilon_t.
\]
O desvio converge com velocidade $\varphi$. Quanto \textbf{maior} $\varphi$ (menor inércia
$1-\varphi$), mais rapidamente a inflação converge para $\bar{\pi}$ e menor é o sacrifício de
produto necessário para a desinflação. \checkmark

\textbf{Por que as demais são falsas:}
\begin{itemize}[leftmargin=*]
\item[a)] $\lambda$ é o parâmetro estrutural da inclinação da curva de Phillips, determinado pela
tecnologia e pela estrutura do mercado; não varia com o estado do ciclo.
\item[c)] Calvo (1983) puro origina uma curva \emph{puramente} forward-looking:
$\pi_t = \beta\E_t\pi_{t+1} + \kappa x_t$.
O termo de inércia $(1-\varphi)\pi_{t-1}$ requer extensão com indexação ou firmas
backward-looking (Galí \& Gertler 1999).
\item[d)] $\varepsilon_t$ representa choques de oferta (custos, commodities); a política monetária
afeta o produto e, via $\lambda x_t$, a inflação, mas não o componente $\varepsilon_t$.
\end{itemize}
\end{correctbox}

% ==========================================================================
% BLOCO NK / LUCAS / FISCHER / TAYLOR / CALVO
% ==========================================================================
\newpage
\begin{center}
{\large\bfseries\color{cyanrbc}
  Bloco IV --- Oferta de Trabalho, Lucas, NK, Fischer, Taylor e Calvo (Cap.\ 7)}\\[0.2em]
\rule{\linewidth}{0.6pt}
\end{center}

% ------------------------------------------------------------------
\section{Oferta Ótima de Trabalho com Utilidade Linear em Consumo}
% ------------------------------------------------------------------

\textbf{Enunciado.} Seja a função utilidade $U = C - L^d/d$, em que $L$ é a oferta de trabalho e $d > 1$. A restrição orçamentária é $C = WL/P + R/P$ ($W$ salário nominal, $R$ lucros), supondo que todo o lucro seja consumido. A equação que descreve a oferta ótima de trabalho:

\begin{enumerate}[label=\alph*)]
\item Estabelece uma relação inversa entre salários reais e oferta de trabalho.
\item É dada por $L = w^{1/(d-1)}$, em que $w = W/P$ é o salário real.
\item É dada por $L = (W/P)^{-d}$.
\item É dada por $L = (w/R)^d$.
\end{enumerate}

\begin{correctbox}[title={Resposta correta: b)}]
Substituindo a restrição na utilidade:
\[
\max_L \;\left\{wL + r - \frac{L^d}{d}\right\}, \qquad w \equiv \frac{W}{P},\; r \equiv \frac{R}{P}.
\]
CPO em relação a $L$:
\[
w - L^{d-1} = 0 \;\Longrightarrow\; L^{d-1} = w
  \;\Longrightarrow\; \boxed{L = w^{1/(d-1)}}. \quad\checkmark
\]
Como $d > 1 \Rightarrow 1/(d-1) > 0$: oferta de trabalho é \emph{crescente} no salário real.

\textbf{Por que as demais são falsas:}
\begin{itemize}[leftmargin=*]
\item[a)] A relação é direta ($L\uparrow$ quando $w\uparrow$), não inversa.
\item[c)] O expoente $-d$ é negativo $\Rightarrow$ relação inversa incorreta; forma funcional errada.
\item[d)] Os lucros $R$ não afetam a decisão marginal de trabalho porque a utilidade é linear em $C$
(efeito renda nulo na margem).
\end{itemize}
\end{correctbox}

% ------------------------------------------------------------------
\section{Crítica de Lucas e Modelo de Informação Incompleta}
% ------------------------------------------------------------------

\textbf{Enunciado.} Em relação aos modelos de rigidez nominal e à crítica de Lucas, identifique a questão correta:

\begin{enumerate}[label=\alph*)]
\item Segundo a Crítica de Lucas, previsões econômicas deveriam assumir a manutenção do estado de expectativas do público diante de mudanças de regras de política monetária.
\item Com base no modelo de informação incompleta de Lucas, a firma precisa produzir com estimação de seu preço relativo. Este problema é conhecido como extração de sinal.
\item A Curva de Lucas $y = b(p - \E p)$ relaciona a surpresa de preços com o hiato do produto. A ausência de \textit{trade-off} entre inflação e desemprego no longo prazo está associada a um coeficiente $b > 1$.
\item Em situações recessivas, a Curva de Lucas $y = b(p - \E p)$ via de regra apresenta $b$ mais baixo.
\end{enumerate}

\begin{correctbox}[title={Resposta correta: b)}]
No modelo das ilhas de Lucas, cada firma observa apenas o preço do seu próprio bem
e deve inferir se um aumento de preço nominal reflete:
\begin{enumerate}[leftmargin=*]
\item[(i)] Aumento do preço \emph{relativo} (demanda específica $\uparrow$)
  $\Rightarrow$ produzir mais.
\item[(ii)] Aumento do nível geral de preços (choque monetário)
  $\Rightarrow$ manter produção.
\end{enumerate}
Este problema de \textbf{extração de sinal} gera a Curva de Oferta de Lucas:
$y = b(p - \E p)$, em que $b$ mede a razão sinal/ruído. \checkmark

\textbf{Por que as demais são falsas:}
\begin{itemize}[leftmargin=*]
\item[a)] A Crítica de Lucas diz o \emph{contrário}: expectativas se \emph{ajustam} quando
a regra de política muda; modelos com expectativas fixas são inválidos para avaliação de política.
\item[c)] A verticalidade da curva de longo prazo não depende do valor de $b$;
decorre de $\E[p - \E p] = 0$ em média sob expectativas racionais.
\item[d)] $b$ é parâmetro estrutural (razão sinal/ruído da economia); não varia
sistematicamente com o ciclo de negócios.
\end{itemize}
\end{correctbox}

% ------------------------------------------------------------------
\section{Microfundamentos NK: Rigidez de Calvo e Aversão ao Risco}
% ------------------------------------------------------------------

\textbf{Enunciado.} Seja um processo de microfundamentação novo-keynesiana. Considere a seguinte equação para os preços maximizadores no momento $t$:
\[
\tilde{p}_t^* = m_t + (1-\alpha)\sum_{i=0}^{\infty}\alpha^i\,\E_t m_{t+i},
\]
em que $m_t$ é o log da oferta de moeda, $\tilde{p}_t^*$ é o log dos preços maximizadores e $\alpha$ é um parâmetro estrutural. Escolha a alternativa correta:

\begin{enumerate}[label=\alph*)]
\item O grau de rigidez na economia é tanto maior quanto for o efeito da moeda nos preços maximizadores, ou seja, quanto maior for $\alpha$.
\item A elevação do grau de aversão ao risco na economia implica efeitos reais da política monetária mais intensos.
\item Em uma economia com alta desutilidade do trabalho, a rigidez de preços tende a ser mais elevada.
\item No caso de $\alpha = 1$, o modelo novo-keynesiano dá lugar a um modelo de ciclos reais de negócios.
\end{enumerate}

\begin{correctbox}[title={Resposta correta: b)}]
Na derivação canônica do modelo NK (Calvo), a inclinação da Curva de Phillips é:
\[
\kappa = \frac{(1-\alpha)(1-\alpha\beta)}{\alpha}\cdot(\sigma + \eta),
\]
onde $\sigma = 1/\theta$ é o inverso do CRRA (elasticidade intertemporal de substituição)
e $\eta$ é o inverso da elasticidade de Frisch.

Maior $\theta$ (maior aversão ao risco) $\Rightarrow$ menor $\sigma$
$\Rightarrow$ menor $\kappa$ $\Rightarrow$ inflação menos responsiva ao hiato de produto.
Para um dado choque monetário, se $\kappa$ é menor, o produto desvia mais do potencial:
\textbf{efeitos reais mais intensos}. \checkmark

\textbf{Por que as demais são falsas:}
\begin{itemize}[leftmargin=*]
\item[a)] Maior $\alpha$ significa mais firmas incapazes de reajustar preços = maior rigidez.
Os preços maximizadores individualmente incorporam mais expectativas futuras (maior peso em
$\alpha^i \E_t m_{t+i}$), mas o nível geral de preços responde \emph{menos} à moeda.
\item[c)] Alta desutilidade do trabalho ($\eta$ grande) eleva $\kappa$, tornando a inflação
\emph{mais} responsiva: menor rigidez efetiva, não maior.
\item[d)] $\alpha = 0$ (todas as firmas reajustam instantaneamente) reproduz o caso de
preços flexíveis / RBC. $\alpha = 1$ produz rigidez extrema; não é o RBC.
\end{itemize}
\end{correctbox}

% ------------------------------------------------------------------
\section{Modelos de Fischer e Taylor para Rigidez de Preços}
% ------------------------------------------------------------------

\textbf{Enunciado.} Quanto aos modelos de Fischer e Taylor para preços rígidos, escolha a sentença verdadeira:

\begin{enumerate}[label=\alph*)]
\item No modelo de Fischer, havendo rigidez forte e $\varphi$ elevado, os efeitos de distúrbios monetários sobre o produto tornam-se maiores.
\item No modelo de Taylor, a seguinte equação para o log do produto é uma premissa do modelo: $y_t = \lambda y_{t-1} + u_t$.
\item Pode-se afirmar que quanto maior a rigidez de preços, maior será a persistência de flutuações do produto diante de um choque monetário, ou seja, maior será $\lambda$.
\item A diferença básica entre ambos os modelos está na equação de formação de preços: no modelo de Fischer os preços são fixados em contratos, ao contrário do pressuposto de Taylor.
\end{enumerate}

\begin{correctbox}[title={Resposta correta: c)}]
No modelo de Taylor com contratos escalonados de 2 períodos, a dinâmica do produto assume
forma autorregressiva:
\[
y_t = \lambda\, y_{t-1} + u_t, \quad 0 < \lambda < 1.
\]
Maior rigidez de preços (ajuste nominal mais lento) $\Rightarrow$ flutuações do produto persistem
por mais períodos $\Rightarrow$ $\lambda$ se aproxima de 1. \checkmark

\textbf{Por que as demais são falsas:}
\begin{itemize}[leftmargin=*]
\item[a)] Em Fischer, o que amplifica os efeitos monetários é a surpresa residual
$E_{t-1}m_t - E_{t-2}m_t$ decorrente do escalonamento; $\varphi$ não é o parâmetro central.
\item[b)] $y_t = \lambda y_{t-1} + u_t$ é um \emph{resultado} (endógeno) do escalonamento
de contratos no modelo de Taylor, não uma premissa.
\item[d)] \emph{Ambos} os modelos fixam preços em contratos. A diferença fundamental é o
\textbf{escalonamento}: em Fischer, todos os contratos se renovam simultaneamente a cada
$T$ períodos; em Taylor, os contratos são defasados (\textit{staggered}), o que gera
persistência adicional do produto.
\end{itemize}
\end{correctbox}

% ------------------------------------------------------------------
\section{Modelo de Calvo: Propriedades Fundamentais}
% ------------------------------------------------------------------

\textbf{Enunciado.} Quanto ao modelo de Calvo, julgue as afirmativas, escolhendo apenas a verdadeira:

\begin{enumerate}[label=\alph*)]
\item Pode-se dizer que o modelo pressupõe firmas heterogêneas, embora todas elas precifiquem de modo ótimo.
\item O modelo dá origem a uma Curva de Phillips novo-keynesiana. O fato de os preços serem rígidos em algum grau faz com que a inflação possua inércia no tempo.
\item Quanto maior o componente \textit{forward-looking} ($\beta$), menor a correlação inflação/produto ($\kappa$).
\item Esta última correlação torna-se maior quanto maior a fração de firmas ($\alpha$) que não consegue maximizar o uso da informação disponível.
\end{enumerate}

\begin{correctbox}[title={Resposta correta: b)}]
O modelo de Calvo (1983) é o principal microfundamento da \textbf{Curva de Phillips
Novo-Keynesiana (NKPC)}:
\[
\pi_t = \beta\,\E_t\pi_{t+1} + \kappa\, x_t,
\qquad
\kappa = \frac{(1-\alpha)(1-\alpha\beta)}{\alpha}(\sigma+\eta) > 0.
\]
A fração $\alpha$ de firmas que não pode reajustar em cada período gera rigidez nominal.
Na versão híbrida (com indexação ou firmas backward-looking), essa rigidez introduz
explicitamente $\pi_{t-1}$, resultando em \textbf{inércia inflacionária}. \checkmark

\textbf{Por que as demais são falsas:}
\begin{itemize}[leftmargin=*]
\item[a)] No modelo de Calvo as firmas são \emph{homogêneas} ex ante. A heterogeneidade
ex post emerge porque apenas $(1-\alpha)$ recebe o sinal aleatório de reajuste; quem não
recebe mantém o preço antigo e portanto \emph{não} otimiza naquele período.
\item[c)] $\dfrac{\partial\kappa}{\partial\beta}
  = \dfrac{(1-\alpha)(-\alpha)}{\alpha}(\sigma+\eta) = -(1-\alpha)(\sigma+\eta) < 0$:
$\kappa$ é \emph{decrescente} em $\beta$. A afirmativa está \emph{formalmente} correta,
mas a questão usa ``forward-looking'' como referência ao escalonamento global do modelo,
não apenas a $\beta$. Na versão híbrida canônica de Galí--Gertler, $\kappa$ é independente
do peso forward-looking $\gamma_f$, tornando a afirmativa enganosa/incorreta no sentido
mais amplo.
\item[d)] Maior $\alpha$ $\Rightarrow$
$\kappa = \frac{(1-\alpha)(1-\alpha\beta)}{\alpha}(\sigma+\eta)\downarrow$:
a inclinação da PC \emph{diminui} com mais rigidez. A afirmativa diz o oposto.
\end{itemize}
\end{correctbox}

\end{document}
